% Part: sets-functions-relations
% Chapter: relations
% Section: special-properties

\documentclass[../../../include/open-logic-section]{subfiles}

\begin{document}

\olfileid{sfr}{rel}{prp}
\olsection{Eigenschappen die vaak voorkomen}

\begin{intro}
Enkele soorten relaties komen zo vaak voor dat ze een eigen naam
hebben. Neem $\le$ en~$\subseteq$. De ene werkt bijvoorbeeld op $\Nat$
(voor~$\le$) en de andere op $\Pow{A}$ (voor~$\subseteq$), maar de
elementen in die domeinen staan op een vergelijkbare manier tot elkaar.
Om precies te zien wat hetzelfde is en wat anders, kijken we naar een
paar eigenschappen van relaties.
Sommige van die eigenschappen, of combinaties ervan, geven ons zeer
belangrijke soorten relaties: ordeningen en equivalentierelaties.
\end{intro}

\begin{defn}[Reflexiviteit]
We noemen een relatie $R \subseteq A^2$ \emph{reflexief} precies als
de relatie ieder element aan zichzelf koppelt: voor elke
$x \in A$ geldt $Rxx$.
\end{defn}

\begin{defn}[Transitiviteit]
We noemen een relatie $R \subseteq A^2$ \emph{transitief} precies als
dit altijd geldt: als $Rxy$ en $Ryz$, dan ook $Rxz$.
\end{defn}

\begin{defn}[Symmetrie]
We noemen een relatie~$R \subseteq A^2$ \emph{symmetrisch} precies als
de relatie ook in de andere richting geldt: als $Rxy$, dan ook~$Ryx$.
\end{defn}

\begin{defn}[Antisymmetrie]
We noemen een relatie~$R \subseteq A^2$ \emph{an\-ti\-symmetrisch}
precies als twee elementen alleen in beide richtingen met elkaar in
relatie kunnen staan wanneer ze gelijk zijn. Dus: als zowel $Rxy$ als
$Ryx$, dan $x=y$ (of, anders gezegd: als $x\neq y$, dan $\lnot Rxy$ of
$\lnot Ryx$).
\end{defn}

\begin{explain}
Bij een symmetrische relatie gaan $Rxy$ en $Ryx$ altijd samen: ze
gelden allebei of geen van beide. Bij een antisymmetrische relatie
kunnen $Rxy$ en $Ryx$ alleen allebei gelden als $x = y$. Dit \emph{betekent niet}
dat $Rxy$ en $Ryx$ moeten gelden wanneer $x = y$; het is alleen niet
uitgesloten. Een antisymmetrische relatie kan dus reflexief zijn, maar
dat geldt niet voor elke antisymmetrische relatie. Ook is
``antisymmetrisch'' niet gewoon een ander woord voor ``niet
symmetrisch''. Een relatie kan tegelijk symmetrisch en
antisymmetrisch zijn. De identiteitsrelatie is daar een voorbeeld van.
\end{explain}

\begin{defn}[Samenhangendheid]
We noemen een relatie $R \subseteq A^2$ \emph{samenhangend} als de
relatie elk tweetal verschillende elementen in minstens één richting
met elkaar verbindt: voor alle $x,y\in A$ geldt dat als $x \neq y$,
dan $Rxy$ of~$Ryx$.
\end{defn}

\begin{prob}
Bedenk voorbeelden van relaties die (a) reflexief en symmetrisch maar
niet transitief zijn, (b) reflexief en antisymmetrisch zijn, (c)
antisymmetrisch en transitief maar niet reflexief zijn, en (d)
reflexief, symmetrisch en transitief zijn. Neem daarbij geen relaties
tussen getallen of tussen verzamelingen.
\end{prob} 
 
\begin{defn}[Irreflexiviteit]
We noemen een relatie $R \subseteq A^2$ \emph{irreflexief} als de
relatie geen enkel element aan zichzelf koppelt: voor alle
$x \in A$ geldt niet $Rxx$.
\end{defn}

\begin{defn}[Asymmetrie]
We noemen een relatie $R \subseteq A^2$ \emph{asymmetrisch} als geen
enkel paar in beide richtingen met elkaar in relatie staat: voor geen
enkel paar $x,y\in A$ gelden zowel $Rxy$ als~$Ryx$.
\end{defn}

Als $A \neq \emptyset$, kan een irreflexieve relatie op~$A$ nooit
reflexief zijn. Elke asymmetrische relatie op~$A$ is wel
antisymmetrisch. Toch zijn er relaties $R \subseteq A^2$ die niet
reflexief maar ook niet irreflexief zijn. En er zijn antisymmetrische
relaties die niet asymmetrisch zijn.
\end{document}